% Generated by scripts/2026-10-21/extract_floats.py from paper/main_ko.tex; edit the full manuscript, not this file.
\begin{table}[htbp]
\centering\small
\caption{Exploratory: linear-probe AUROC by reader and depth, first 250 episodes, mean over steps 0--2.}\label{tab:a6}
\begin{tabularx}{\linewidth}{@{}lccccL@{}}
\toprule
Reader & 25\% & 50\% & 75\% & 100\% & Probe $-$ rung 5, all 24 cells per reader (min, max) \\
\midrule
0.5B & 0.618 & 0.615 & \textbf{0.660} & 0.642 & $-0.277$, $+0.001$ \\
1.5B & 0.589 & 0.631 & \textbf{0.717} & 0.675 & $-0.195$, $-0.006$ \\
3B (= actor) & 0.595 & 0.594 & \textbf{0.727} & 0.682 & $-0.205$, $+0.004$ \\
7B (4-bit) & 0.608 & 0.643 & 0.704 & \textbf{0.717} & $-0.191$, $-0.004$ \\
\bottomrule
\end{tabularx}
\par\smallskip\raggedright\footnotesize
5단은 step 0--2 평균 0.752였다. 마지막 열은 네 깊이·세 step의 선형·MLP 프로브를 모두 포함한다(reader당 24칸, 모두 96칸). 굵은 글씨는 reader별 최고 깊이이다. 3B reader 50\% 깊이가 사전 등록 조건이다(표~\ref{tab:ladder}). 96칸 모두의 차이와 95\% CI는 표~\ref{tab:a9}에 있다.
% 출처: results/explore_{0.5B,1.5B,3B,7B-4bit}.csv, *_diff.csv
\end{table}
